% ============================================================
% CAPITOLUL 9 — APLICAȚIA WEB E-MEDICINE
% ============================================================

\chapter{Aplicația Web E-Medicine}
\label{ch:aplicatie}

\section{Rolul aplicației}
\label{sec:functionalitati}

Aplicația web este interfața prin care modelul matematic și infrastructura de
execuție devin utilizabile. Ea nu este un sistem clinic complet, ci un mediu de
lucru pentru organizarea datelor pe studii epidemiologice, încărcarea
dataseturilor, configurarea simulărilor SIDARTHE, SIR, SIS, SEIR sau SIRS și
compararea rezultatelor obținute în R, Julia și Python. Din această
perspectivă, aplicația are rol de instrument experimental și colaborativ.

Fluxul principal începe cu autentificarea utilizatorului. După autentificare,
acesta alege un studiu sau creează unul nou. Toate operațiile ulterioare se
raportează la studiul curent, iar interfața afișează acțiunile disponibile în
funcție de rolul utilizatorului: \texttt{admin}, \texttt{researcher},
\texttt{contributor} sau \texttt{viewer}. Utilizatorii cu rol suficient pot
încărca un fișier CSV cu date epidemiologice, pot alege modelul, pot introduce
parametrii și pot urmări statusul execuției. Când runner-ul finalizează job-ul,
interfața afișează traiectoriile compartimentale și permite exportul
rezultatelor în format CSV. O pagină separată permite compararea mai multor
rulări finalizate, iar pagina de benchmark prezintă agregările calculate la
nivel de model și limbaj.

Pe lângă aplicația principală există o interfață \texttt{frontend-admin},
servită de pe host-ul \texttt{admin.sir-workbench.ro}. Aceasta este disponibilă doar
conturilor cu \texttt{is\_super\_admin=true} și oferă statistici globale,
gestionarea utilizatorilor și vizualizarea tuturor studiilor.

\section{Backend Go}
\label{sec:backend_go}

Backend-ul este implementat în Go și expune un API REST sub prefixul
\texttt{/api}. Serverul folosește routerul \texttt{chi}, middleware pentru
logging, recuperare după panică, timeout de request și CORS, iar accesul la rutele
protejate este controlat prin JWT. Alegerea Go este potrivită pentru această
componentă deoarece API-ul este stateless, are nevoie de latențe predictibile și
poate fi distribuit ușor sub formă de binar static în container.

Structura codului separă intrarea HTTP de accesul la date. Pachetul
\texttt{handlers} parsează cererile, validează input-ul și construiește
răspunsurile JSON. Pachetul \texttt{repository} conține interogările SQL și
operațiile tranzacționale peste \texttt{pgx}. Pachetul \texttt{runner} definește
schema mesajului trimis în Redis și publicarea job-urilor. Contractul de simulare
include tipul modelului, parametrii, condițiile inițiale și conexiunea la baza de
date.

Rutele aplicației principale sunt organizate sub
\texttt{/api/studies/\{studyId\}}. Middleware-ul RBAC verifică membership-ul și
rolul înainte ca handler-ul să ajungă la date. Super-adminii au acces global,
dar consola administrativă folosește rute dedicate sub \texttt{/api/admin}.

\section{Gestionarea dataseturilor}
\label{sec:dataseturi_web}

Dataseturile sunt încărcate printr-un formular multipart, în contextul studiului
curent. API-ul acceptă câmpurile \texttt{name}, \texttt{disease\_code} și
\texttt{file}; fișierul trebuie să conțină cel puțin coloanele \texttt{date},
\texttt{region}, \texttt{infected} și \texttt{population}. Coloanele
\texttt{recovered} și \texttt{susceptible} sunt opționale, dar atunci când sunt
prezente sunt folosite pentru condițiile inițiale. Pentru SIDARTHE, cazurile
active sunt plasate inițial în compartimentul $I$, persoanele ieșite din
infecție în $H$, iar compartimentul $E$ al deceselor pornește de la zero.
Pentru SEIR, compartimentul expușilor pornește de asemenea de la zero,
deoarece nu este observat direct în CSV.

Persistarea datasetului se face într-o tranzacție. Mai întâi este creat rândul din
\texttt{datasets}, apoi înregistrările epidemiologice sunt inserate prin
\texttt{CopyFrom}. Această alegere reduce costul încărcării CSV-urilor față de
inserarea rând-cu-rând și păstrează atomicitatea: dacă una dintre înregistrări nu
poate fi scrisă, setul de date nu rămâne într-o stare parțială. La ștergere,
\texttt{admin}-ul poate elimina orice set de date din studiu, în timp ce
\texttt{contributor}-ul poate elimina doar seturile de date încărcate de el.

\section{Execuția simulărilor}
\label{sec:executie_simulari_web}

Crearea unei simulări este un proces scurt în API și un proces asincron în
ansamblul sistemului. Endpointul
\texttt{POST /api/studies/\{studyId\}/simulations} verifică existența
datasetului în studiul curent, modelul, limbajul selectat, parametrii și
$t_{\max}$, apoi extrage prima înregistrare epidemiologică pentru a determina
valorile inițiale. După ce rândul din \texttt{simulation\_runs} este creat cu
status \texttt{pending}, API-ul publică mesajul JSON în coada Redis
corespunzătoare limbajului.

Runner-ele sunt responsabile pentru tranzițiile ulterioare. Ele marchează rularea
ca \texttt{running}, execută solver-ul pentru \texttt{model\_type}, inserează
rezultatele și benchmark-ul, apoi marchează rularea ca \texttt{completed}. Dacă
apare o eroare, statusul devine \texttt{failed}, iar mesajul este salvat în baza de
date. Interfața web nu trebuie să țină o conexiune deschisă pe durata calculului;
pagina de rezultat interoghează periodic endpointul de status până la finalizare.

\section{Frontend Vue.js}
\label{sec:frontend_vue}

Frontend-ul principal este o aplicație Vue 3 construită cu Vite și TypeScript.
Rutarea este gestionată cu Vue Router, iar starea aplicației cu Pinia. Clientul
HTTP este o instanță Axios configurată cu prefixul \texttt{/api}; un interceptor
adaugă automat headerul \texttt{Authorization} atunci când există token în
\texttt{localStorage}. Store-ul \texttt{studies.store.ts} păstrează studiul
curent și rolul curent, iar store-urile pentru seturi de date, simulări și
benchmark-uri construiesc URL-urile în funcție de acest \texttt{studyId}.

Paginile sunt organizate după fluxurile principale ale aplicației:
\texttt{StudiesPage} selectează contextul de lucru,
\texttt{CreateStudyPage} creează un studiu nou,
\texttt{StudyManagementPage} gestionează membrii și rolurile,
\texttt{DashboardPage} sumarizează seturile de date și simulările recente,
\texttt{DatasetsPage} gestionează încărcarea CSV-urilor,
\texttt{NewSimulationPage} colectează modelul și parametrii,
\texttt{SimulationResultPage} afișează statusul și rezultatele,
\texttt{ComparePage} suprapune seriile mai multor rulări, iar
\texttt{BenchmarksPage} afișează agregările pe model și limbaj. Componenta
\texttt{CompartmentChart} folosește Chart.js prin \texttt{vue-chartjs} și
construiește dinamic seriile în funcție de cheile disponibile în
\texttt{compartments}. Astfel, același grafic poate afișa modele cu două,
trei, patru sau opt compartimente fără modificări suplimentare.

Decizia de a folosi polling la două secunde pentru actualizarea statusului unei
simulări este simplă, dar adecvată pentru durata așteptată a job-urilor. O
soluție bazată pe WebSocket sau Server-Sent Events ar reduce numărul de cereri,
însă ar introduce stare suplimentară în backend. Pentru un API stateless și
pentru job-uri care durează de la câteva secunde la câteva minute, polling-ul
păstrează implementarea previzibilă.

\section{Consola super-admin}
\label{sec:frontend_admin}

Consola administrativă este un build Vue separat, cu propriul router și propriul
token în \texttt{localStorage}. La autentificare, aplicația verifică răspunsul
\texttt{/api/auth/login}; dacă utilizatorul nu este super-admin, tokenul nu este
păstrat și login-ul este refuzat. Această separare reduce riscul ca acțiunile
globale de administrare să fie amestecate cu fluxurile obișnuite de lucru ale
unui studiu. Toate paginile sunt deliberat \emph{de citire}: mutațiile asupra
studiilor și seturilor de date rămân în aplicația principală, controlate prin
RBAC pe studiu.

\textbf{Dashboard.} Pagina inițială afișează cele opt contoare globale
(\texttt{/api/admin/stats}) și șase distribuții agregate între studii, obținute
din \texttt{/api/admin/breakdown}: simulări pe status, model și limbaj;
seturi de date și studii pe boală; top zece studii după volumul cumulat de
înregistrări. Distribuțiile sunt randate ca bare orizontale proporționale, iar
top-ul de studii are link-uri direct către pagina de detaliu.

\textbf{Utilizatori și studii.} Pagina de utilizatori
(\texttt{/api/admin/users}) listează toate conturile cu un buton care comută
flag-ul \texttt{is\_super\_admin} prin \texttt{PATCH
/api/admin/users/\{userId\}/super-admin}. Pagina de studii
(\texttt{/api/admin/studies}) listează toate studiile cu un câmp de filtrare
text și permite navigarea către pagina de detaliu prin click pe rând.

\textbf{Detaliu studiu.} \texttt{StudyDetailPage} reutilizează endpointurile
\texttt{/api/studies/\{studyId\}/*} (membri, seturi de date, simulări,
benchmark-uri agregate), de care super-adminul beneficiază datorită
bypass-ului din middleware-ul RBAC. Pagina prezintă, în ordine, un antet cu
metadatele studiului și patru contoare (membri, seturi de date cu volum
cumulat, simulări, combinații model$\,\times\,$limbaj cu rulări complete),
urmate de tabelele detaliate. Simulările au badge de status colorat, iar
tabelul de benchmark-uri arată wall-time median, percentila 95, CPU median și
memorie mediană pentru fiecare combinație model--limbaj.

\textbf{Activitate.} Pagina \texttt{/activity} consumă \texttt{/api/admin/recent-simulations}
și \texttt{/api/admin/recent-datasets} pentru a afișa cele mai recente $N$
intrări (configurabil între 20 și 200), indiferent de studiu, cu link-uri către studiul
corespunzător. Acest feed este util pentru audit operațional, pentru
identificarea rapidă a simulărilor blocate în starea \texttt{running} și pentru
verificarea încărcărilor de seturi de date proaspete.

\section{Validare funcțională}
\label{sec:testare_functionala}

Validarea funcțională a aplicației poate fi făcută prin fluxuri end-to-end care
urmează operațiile reale ale utilizatorului: înregistrare, autentificare, creare
sau selectare studiu, adăugare membru, încărcare set de date, pornire simulare
SIDARTHE, așteptare finalizare, afișare grafic și export CSV. Aceeași suită ar
trebui să verifice și rulările SIR/SIS/SEIR/SIRS, deoarece acestea asigură
compatibilitatea cu familia de modele discutată în capitolul teoretic și cu
comparațiile simple. Un flux separat validează consola
admin: login cu un cont super-admin, listarea utilizatorilor, modificarea
flag-ului \texttt{is\_super\_admin} și listarea studiilor globale.
